\chapter{First Appendix}
\label{app:first}

This is an example of an appendix chapter. Appendices automatically use alphabetic lettering (Appendix A, Appendix B) instead of standard numeric chapters.

\section{Supplementary Information}
You can include extra derivations, long tables of data, or additional code snippets here that are relevant but too detailed for the main text.

\Cref{tab:supp_calibration} lists supplementary calibration parameters at the beginning of a sentence. Notice how opening a sentence with \mintinline{latex}{\Cref} produces the capitalized form (\Cref{tab:supp_calibration}), while mid-sentence references such as \cref{tab:supp_calibration} preserve lowercase, with both adopting the appendix letter prefix \mintinline{text}{A.1}.

\begin{table}[htpb]
    \centering
    \caption{Supplementary detector calibration parameters.}
    \label{tab:supp_calibration}
    \begin{tabular}{lccc}
        \toprule
        Receiver Channel & Center Frequency (MHz) & Gain (dB) & Noise Temp (K) \\
        \midrule
        CH1 (L-Band) & 1420.4 & 45.2 & 32.1 \\
        CH2 (OH-Line) & 1665.4 & 44.8 & 34.5 \\
        CH3 (OH-Line) & 1667.3 & 45.0 & 33.8 \\
        \bottomrule
    \end{tabular}
\end{table}
